% ===== PRÉAMBULE CANONIQUE — à coller VERBATIM en tête de CHAQUE .tex =====
\documentclass[11pt,a4paper]{article}
\usepackage[utf8]{inputenc}\usepackage[T1]{fontenc}\usepackage{lmodern}
\usepackage[french]{babel}
\usepackage{amsmath,amssymb,mathtools}\usepackage{icomma}
\usepackage{siunitx}\sisetup{locale=FR}  % \qty{}{}, \si{}, \ang{}
\usepackage{xcolor}\usepackage{colortbl}
\usepackage{tikz}\usepackage{pgfplots}\pgfplotsset{compat=1.18}
\usepackage[siunitx,europeanresistors]{circuitikz} % circuits (élec.)
\usepackage[most]{tcolorbox}\usepackage{enumitem}\usepackage{array,booktabs}
\usepackage[a4paper,margin=2.2cm]{geometry}
\usetikzlibrary{arrows.meta,calc,angles,quotes,positioning,patterns,%
  backgrounds,shapes.geometric,decorations.pathmorphing,decorations.markings,intersections}
\definecolor{primary}{RGB}{222,49,99}\definecolor{primarydark}{RGB}{150,22,55}
\definecolor{defcol}{RGB}{0,114,178}\definecolor{methcol}{RGB}{52,92,148}
\definecolor{excol}{RGB}{0,135,90}\definecolor{attncol}{RGB}{200,40,55}
\tcbset{boxbase/.style={enhanced,breakable,boxrule=0.9pt,arc=2.5pt,
  left=3mm,right=3mm,top=2.3mm,bottom=2.3mm,fonttitle=\bfseries,coltitle=white}}
% --- boîtes TUTORIEL (noms EXACTS, ne pas renommer) ---
\newtcolorbox{cle}[1][]{boxbase,colback=defcol!6,colframe=defcol,
  title={\textsc{À retenir}\ifx\relax#1\relax\relax\else\textmd{ : \itshape #1}\fi}}
\newtcolorbox{methode}[1][]{boxbase,colback=methcol!7,colframe=methcol,sharp corners,
  title={\textsc{Méthode}\ifx\relax#1\relax\relax\else\textmd{ --- \itshape #1}\fi}}
\newtcolorbox{exemplebox}[1][]{boxbase,colback=excol!5,colframe=excol,
  title={\textsc{Exemple}\ifx\relax#1\relax\relax\else\textmd{ : \itshape #1}\fi}}
\newtcolorbox{piege}[1][]{boxbase,colback=attncol!6,colframe=attncol,
  title={\textsc{Piège}\ifx\relax#1\relax\relax\else\textmd{ : \itshape #1}\fi}}
\newtcolorbox{remarque}[1][]{boxbase,colback=black!4,colframe=black!45,
  title={\textsc{Remarque}\ifx\relax#1\relax\relax\else\textmd{ : \itshape #1}\fi}}
% --- boîtes EXERCICE / CORRIGÉ (noms EXACTS, auto-comptées) ---
\newtcolorbox[auto counter]{exo}[1][]{boxbase,colback=primary!4,colframe=primary,
  title={\textsc{Exercice}~\thetcbcounter\ifx\relax#1\relax\relax\else\textmd{ --- \itshape #1}\fi}}
\newtcolorbox[auto counter]{corrige}[1][]{boxbase,colback=excol!4,colframe=excol,
  title={\textsc{Corrigé}~\thetcbcounter\ifx\relax#1\relax\relax\else\textmd{ --- \itshape #1}\fi}}
\newcommand{\vv}[1]{\vec{#1}}\newcommand{\ds}{\displaystyle}
\newcommand{\R}{\mathbb{R}}\newcommand{\N}{\mathbb{N}}\newcommand{\Z}{\mathbb{Z}}
% (un \usepackage{fancyhdr}+en-tête est possible ; il est ignoré côté web)
% =========================================================================
\begin{document}
\sisetup{per-mode=symbol}

\begin{corrige}[différence de potentiel]
\begin{enumerate}[label=\alph*)]
  \item $U_{AB} = V_A - V_B = 12 - (-3) = \SI{15}{\volt}$.
  \item $U_{BA} = V_B - V_A = -3 - 12 = \SI{-15}{\volt} = -U_{AB}$.
\end{enumerate}
\end{corrige}

\begin{corrige}[travail de la force électrique]
\begin{enumerate}[label=\alph*)]
  \item $W = q\,U_{AB} = \num{2e-6}\times 100 = \SI{2e-4}{\joule}$ ($>0$).
  \item $W = q\,U_{AB} = (\num{-2e-6})\times 100 = \SI{-2e-4}{\joule}$ ($<0$ : il faut fournir un
        travail contre la force).
\end{enumerate}
\end{corrige}

\begin{corrige}[à quoi correspond un volt ?]
\begin{enumerate}[label=\alph*)]
  \item $U = W/q$ : une énergie divisée par une charge, donc $\SI{1}{\volt} = \SI{1}{\joule\per\coulomb}$.
  \item En remplaçant $\si{\joule}$ par $\si{\newton\metre}$ :
        $\SI{1}{\volt} = \SI{1}{\newton\metre\per\coulomb}$.
\end{enumerate}
\end{corrige}

\begin{corrige}[champ dans un condensateur]
\begin{enumerate}[label=\alph*)]
  \item $E = \dfrac{|U_{AB}|}{d} = \dfrac{60}{0{,}03} = \SI{2000}{\volt\per\metre}$.
  \item $E = \SI{2000}{\newton\per\coulomb}$ (les deux unités sont identiques).
\end{enumerate}
\end{corrige}

\begin{corrige}[le travail double avec la charge]
\begin{enumerate}[label=\alph*)]
  \item Non : la différence de potentiel ne dépend que des charges \emph{sources} et des positions,
        pas de la charge déplacée. $U_{AB}$ est donc la même dans les deux cas.
  \item $W = q\,U_{AB}$ est proportionnel à la charge : en la multipliant par $10$, le travail est
        multiplié par $10$ :
        \[ W' = 10\,Q\,U_{AB} = 10\,x\ \text{joules}. \]
\end{enumerate}
\end{corrige}

\end{document}
